
\documentclass[12pt,a4paper]{article}

% ----- Pacotes -----
\usepackage[utf8]{inputenc}
\usepackage[T1]{fontenc}
\usepackage[brazilian]{babel}
\usepackage{geometry}
\usepackage{graphicx}
\usepackage{listings}
\usepackage{xcolor}
\usepackage{hyperref}
\usepackage{enumitem}
\usepackage{titlesec}
\usepackage{fancyhdr}
\usepackage{mdframed}
\usepackage{booktabs}
\usepackage{amsmath}
\usepackage{tcolorbox}
\usepackage{array}
\tcbuselibrary{listingsutf8, skins, breakable}

\geometry{top=2.5cm, headheight=15pt, bottom=2.5cm, left=3cm, right=2.5cm}

% ----- Cores -----
\definecolor{codeblue}{RGB}{0,100,180}
\definecolor{codegray}{RGB}{245,245,245}
\definecolor{codegreen}{RGB}{0,128,0}
\definecolor{codered}{RGB}{180,0,0}
\definecolor{alertyellow}{RGB}{255,248,220}
\definecolor{alertborder}{RGB}{200,160,0}
\definecolor{tipgreen}{RGB}{230,255,230}
\definecolor{tipborder}{RGB}{0,150,0}
\definecolor{headerblue}{RGB}{30,80,150}
\definecolor{dockerblue}{RGB}{0,100,150}

% ----- Estilo terminal -----
\lstdefinestyle{terminal}{
  backgroundcolor=\color{black!90},
  basicstyle=\ttfamily\small\color{white},
  breaklines=true,
  frame=single,
  numbers=none,
  xleftmargin=0pt,
  framexleftmargin=0pt
}

% ----- Estilo Python -----
\lstdefinestyle{python}{
  backgroundcolor=\color{codegray},
  basicstyle=\ttfamily\small,
  keywordstyle=\color{codeblue}\bfseries,
  commentstyle=\color{codegreen}\itshape,
  stringstyle=\color{codered},
  breaklines=true,
  frame=single,
  numbers=left,
  numberstyle=\tiny\color{gray},
  xleftmargin=1.5em,
  framexleftmargin=1.5em,
  language=Python,
  showstringspaces=false,
  literate={a}{{a}}1
}

% ----- Estilo código genérico -----
\lstdefinestyle{code}{
  backgroundcolor=\color{codegray},
  basicstyle=\ttfamily\small,
  breaklines=true,
  frame=single,
  numbers=left,
  numberstyle=\tiny\color{gray},
  xleftmargin=1.5em,
  framexleftmargin=1.5em,
  showstringspaces=false
}

% ----- Caixas de destaque -----
\newtcolorbox{alertbox}{
  colback=alertyellow,
  colframe=alertborder,
  fonttitle=\bfseries,
  title={Aten\c{c}\~ao},
  breakable
}


\newtcolorbox{infobox}[1][Informa\c{c}\~ao]{
  colback=blue!5,
  colframe=headerblue,
  fonttitle=\bfseries,
  title={#1},
  breakable
}

\newtcolorbox{questbox}{
  colback=blue!3,
  colframe=dockerblue,
  fonttitle=\bfseries,
  title={Quest\~oes para Reflex\~ao},
  breakable
}

% ----- Comando para campo em branco -----
\newcommand{\campo}[1]{\underline{\hspace{#1}}}

% ----- Cabeçalho e rodapé -----
\pagestyle{fancy}
\fancyhf{}
\fancyhead[L]{\small\textcolor{headerblue}{\textbf{IFCE -- Ci\^encia da Computa\c{c}\~ao}}}
\fancyhead[R]{\small\textcolor{headerblue}{API Gateway com Nginx e Docker Compose}}
\fancyfoot[C]{\thepage}
\renewcommand{\headrulewidth}{0.4pt}

% ----- Formatação de seções -----
\titleformat{\section}
  {\Large\bfseries\color{headerblue}}
  {\thesection.}{0.5em}{}[\titlerule]

\titleformat{\subsection}
  {\large\bfseries\color{headerblue!80}}
  {\thesubsection.}{0.5em}{}

\titleformat{\subsubsection}
  {\normalsize\bfseries\color{headerblue!60}}
  {\thesubsubsection.}{0.5em}{}

% ================================================================
\begin{document}

% ----- Capa -----
\begin{titlepage}
  \centering
  \vspace*{1cm}
  {\Large \textbf{INSTITUTO FEDERAL DE EDUCA\c{C}\~AO, CI\^ENCIA E TECNOLOGIA DO CEAR\'A}}\\[0.3cm]
  {\large Curso: Ci\^encia da Computa\c{c}\~ao}\\[0.1cm]
  {\large Disciplina: \campo{7cm}}\\[0.1cm]
  {\large Professor(a): \campo{6cm}}

  \vspace{2cm}

  \rule{\textwidth}{1.5pt}\\[0.4cm]
  {\Huge \textbf{Roteiro de Aula Pr\'atica}}\\[0.3cm]
  {\LARGE API Gateway com Nginx em Microsservi\c{c}os}\\[0.3cm]
  {\large Docker Compose, Spring Boot e Proxy Reverso}\\[0.4cm]
  \rule{\textwidth}{1.5pt}

  \vspace{2cm}

  \begin{infobox}[Objetivos do Documento]
    Ao final desta aula, o aluno ser\'a capaz de:
    \begin{itemize}[leftmargin=1.5em]
      \item entender por que expor cada microsservi\c{c}o em uma porta diferente pode ser ruim;
      \item explicar o papel de um API Gateway em uma arquitetura distribu\'ida;
      \item configurar o Nginx como proxy reverso para tr\^es APIs REST;
      \item alterar o \texttt{docker-compose.yml} para publicar apenas o gateway;
      \item fazer o front consumir as APIs por uma porta \'unica;
      \item validar a solu\c{c}\~ao com \texttt{docker ps}, \texttt{docker compose config} e \texttt{curl}.
    \end{itemize}
  \end{infobox}

  \vspace{1.2cm}

  \begin{infobox}[Produto Final]
    O produto final deste roteiro \'e uma aplica\c{c}\~ao de loja com:
    \begin{itemize}[leftmargin=1.5em]
      \item \texttt{clientes-service}, \texttt{estoque-service} e \texttt{pedidos-service} acess\'iveis por \texttt{http://localhost:9000};
      \item \texttt{web-ms} mantido em \texttt{http://localhost:9008};
      \item Nginx roteando \texttt{/api/clientes}, \texttt{/api/estoques} e \texttt{/api/pedidos};
      \item portas antigas \texttt{9005}, \texttt{9006} e \texttt{9007} fechadas para acesso externo.
    \end{itemize}
  \end{infobox}

  \vfill
  {\large Fortaleza, \the\year}
\end{titlepage}

\tableofcontents
\newpage

% ================================================================
\section{Como Ler Este Material}

Este roteiro foi pensado para ser usado como guia de uma aula pr\'atica. Ele mostra o problema inicial, a motiva\c{c}\~ao arquitetural, as altera\c{c}\~oes feitas no projeto e os testes usados para confirmar que tudo funcionou.

\begin{infobox}[Estrat\'egia da Aula]
  A aula est\'a dividida em seis partes:
  \begin{enumerate}[leftmargin=2em]
    \item identificar os microsservi\c{c}os e portas abertas;
    \item discutir o problema de expor cada API separadamente;
    \item criar a configura\c{c}\~ao do Nginx;
    \item alterar o \texttt{docker-compose.yml};
    \item subir os containers novamente;
    \item testar o gateway e as portas antigas.
  \end{enumerate}
\end{infobox}

\begin{alertbox}
  O objetivo n\~ao \'e apenas copiar a configura\c{c}\~ao. O ponto principal da aula \'e entender que o gateway se torna a entrada \'unica das APIs, enquanto os microsservi\c{c}os continuam existindo internamente.
\end{alertbox}

% ================================================================
\section{Parte 1 --- Diagn\'ostico Inicial}

\subsection{Contexto}

O projeto \texttt{loja-microservices} possui uma pequena aplica\c{c}\~ao de loja dividida em microsservi\c{c}os. Antes da mudan\c{c}a, cada API estava exposta diretamente em uma porta diferente da m\'aquina host.

\begin{center}
\begin{tabular}{lll}
\toprule
\textbf{Servi\c{c}o} & \textbf{Responsabilidade} & \textbf{Porta antes da mudan\c{c}a} \\
\midrule
\texttt{clientes-service} & API de clientes & \texttt{9005 -> 8081} \\
\texttt{estoque-service} & API de estoque & \texttt{9006 -> 8082} \\
\texttt{pedidos-service} & API de pedidos & \texttt{9007 -> 8083} \\
\texttt{web-ms} & Interface web & \texttt{9008 -> 8080} \\
\bottomrule
\end{tabular}
\end{center}

\subsection{Comandos usados para observar o ambiente}

\begin{lstlisting}[style=terminal]
ss -ltnp
\end{lstlisting}

Esse comando mostra portas TCP em estado de escuta. Ele permitiu observar as portas \texttt{9005}, \texttt{9006}, \texttt{9007} e \texttt{9008} abertas.

\begin{lstlisting}[style=terminal]
docker ps --format 'table {{.Names}}\t{{.Image}}\t{{.Ports}}\t{{.Status}}'
\end{lstlisting}

Esse comando mostrou quais containers estavam rodando, suas imagens, portas publicadas e estado atual.

\begin{infobox}[Resultado observado]
  Os containers principais eram:
  \begin{itemize}[leftmargin=1.5em]
    \item \texttt{web-ms}, publicado em \texttt{9008};
    \item \texttt{clientes-service}, publicado em \texttt{9005};
    \item \texttt{estoque-service}, publicado em \texttt{9006};
    \item \texttt{pedidos-service}, publicado em \texttt{9007};
    \item bancos \texttt{clientes-db}, \texttt{estoque-db} e \texttt{pedidos-db}, acess\'iveis apenas dentro da rede Docker.
  \end{itemize}
\end{infobox}

\subsection{Porta interna e porta externa}

No Docker Compose, uma declara\c{c}\~ao como esta:

\begin{lstlisting}[style=code]
ports:
  - "9005:8081"
\end{lstlisting}

significa que a porta \texttt{9005} da m\'aquina host encaminha para a porta \texttt{8081} dentro do container.


\begin{questbox}
  \begin{itemize}[leftmargin=1.5em]
    \item O consumidor da API deveria conhecer a porta de cada microsservi\c{c}o?
    \item O que acontece se uma API mudar de porta?
    \item Expor v\'arias portas aumenta ou diminui a superf\'icie de acesso da aplica\c{c}\~ao?
  \end{itemize}
\end{questbox}

% ================================================================
\section{Parte 2 --- Problema que Queremos Resolver}

\subsection{Situa\c{c}\~ao antes do gateway}

Antes da altera\c{c}\~ao, para consumir todas as APIs era necess\'ario conhecer v\'arias portas:

\begin{lstlisting}[style=terminal]
http://localhost:9005/api/clientes
http://localhost:9006/api/estoques
http://localhost:9007/api/pedidos
\end{lstlisting}

Isso funciona, mas tem limita\c{c}\~oes importantes:

\begin{itemize}[leftmargin=1.5em]
  \item cada servi\c{c}o fica diretamente exposto no host;
  \item o cliente precisa conhecer detalhes internos da arquitetura;
  \item fica mais dif\'icil centralizar logs, autentica\c{c}\~ao, CORS e limita\c{c}\~ao de requisi\c{c}\~oes;
  \item a arquitetura fica menos parecida com uma solu\c{c}\~ao real de produ\c{c}\~ao.
\end{itemize}

\subsection{Solu\c{c}\~ao proposta}

A solu\c{c}\~ao foi criar um \textbf{API Gateway} com Nginx. Depois disso, as APIs passam a ser acessadas por uma \'unica porta:

\begin{lstlisting}[style=terminal]
http://localhost:9000/api/clientes
http://localhost:9000/api/estoques
http://localhost:9000/api/pedidos
\end{lstlisting}

\begin{infobox}[O que \'e um API Gateway?]
  Um API Gateway \'e um ponto de entrada para chamadas externas. Ele recebe uma requisi\c{c}\~ao, analisa o caminho da URL e encaminha para o servi\c{c}o interno correto. Neste roteiro, o Nginx faz esse papel como proxy reverso.
\end{infobox}

% ================================================================
\section{Parte 3 --- Criando a Configura\c{c}\~ao do Nginx}

\subsection{Criando a pasta}

Foi criada uma pasta chamada \texttt{nginx} na raiz do projeto:

\begin{lstlisting}[style=terminal]
mkdir -p nginx
\end{lstlisting}

\subsection{Criando o arquivo do gateway}

Dentro dessa pasta, foi criado o arquivo:

\begin{lstlisting}[style=terminal]
nginx/api-gateway.conf
\end{lstlisting}

com o seguinte conte\'udo:

\begin{lstlisting}[style=code]
server {
    listen 9000;
    server_name _;

    proxy_http_version 1.1;
    proxy_set_header Host $host;
    proxy_set_header X-Real-IP $remote_addr;
    proxy_set_header X-Forwarded-For $proxy_add_x_forwarded_for;
    proxy_set_header X-Forwarded-Proto $scheme;

    location /api/clientes {
        proxy_pass http://clientes-service:8081;
    }

    location /api/estoques {
        proxy_pass http://estoque-service:8082;
    }

    location /api/pedidos {
        proxy_pass http://pedidos-service:8083;
    }

    location / {
        return 404;
    }
}
\end{lstlisting}

\subsection{Explica\c{c}\~ao da configura\c{c}\~ao}

\begin{itemize}[leftmargin=1.5em]
  \item \texttt{listen 9000}: faz o Nginx escutar na porta \texttt{9000} dentro do container.
  \item \texttt{server\_name \_}: aceita requisi\c{c}\~oes independentemente do nome do host.
  \item \texttt{proxy\_set\_header}: preserva informa\c{c}\~oes da requisi\c{c}\~ao original.
  \item \texttt{location /api/clientes}: define que chamadas para clientes v\~ao para \texttt{clientes-service:8081}.
  \item \texttt{location /api/estoques}: define que chamadas para estoque v\~ao para \texttt{estoque-service:8082}.
  \item \texttt{location /api/pedidos}: define que chamadas para pedidos v\~ao para \texttt{pedidos-service:8083}.
  \item \texttt{location /}: retorna \texttt{404} para qualquer caminho n\~ao configurado.
\end{itemize}


% ================================================================
\section{Parte 4 --- Alterando o docker-compose.yml}

\subsection{Removendo as portas p\'ublicas das APIs}

Antes, os servi\c{c}os de API tinham \texttt{ports}. Exemplo:

\begin{lstlisting}[style=code]
clientes-service:
  build:
    context: ./clientes-service
  container_name: clientes-service
  depends_on:
    - clientes-db
  environment:
    SERVER_PORT: 8081
  ports:
    - "9005:8081"
\end{lstlisting}

A parte removida foi:

\begin{lstlisting}[style=code]
ports:
  - "9005:8081"
\end{lstlisting}

O mesmo foi feito para:

\begin{itemize}[leftmargin=1.5em]
  \item \texttt{estoque-service}, removendo \texttt{9006:8082};
  \item \texttt{pedidos-service}, removendo \texttt{9007:8083}.
\end{itemize}

\begin{alertbox}
  Remover \texttt{ports} n\~ao desliga o servi\c{c}o. Ele continua acess\'ivel por outros containers na mesma rede Docker. A diferen\c{c}a \'e que ele deixa de ficar acess\'ivel diretamente pela m\'aquina host.
\end{alertbox}

\subsection{Adicionando o servi\c{c}o api-gateway}

Foi adicionado este servi\c{c}o ao \texttt{docker-compose.yml}:

\begin{lstlisting}[style=code]
api-gateway:
  image: nginx:1.27-alpine
  container_name: api-gateway
  depends_on:
    - clientes-service
    - estoque-service
    - pedidos-service
  volumes:
    - ./nginx/api-gateway.conf:/etc/nginx/conf.d/default.conf:ro
  ports:
    - "9000:9000"
\end{lstlisting}

\subsection{Detalhes importantes}

\begin{itemize}[leftmargin=1.5em]
  \item \texttt{image: nginx:1.27-alpine}: usa uma imagem oficial e leve do Nginx.
  \item \texttt{container\_name: api-gateway}: define um nome fixo para facilitar comandos e logs.
  \item \texttt{depends\_on}: indica que o gateway depende dos servi\c{c}os de API.
  \item \texttt{volumes}: monta a configura\c{c}\~ao local do Nginx dentro do container.
  \item \texttt{:ro}: monta o arquivo como somente leitura.
  \item \texttt{9000:9000}: publica a porta \texttt{9000} no host.
\end{itemize}

\begin{alertbox}
  \texttt{depends\_on} ajuda na ordem de subida dos containers, mas n\~ao garante que uma aplica\c{c}\~ao Spring Boot j\'a esteja pronta para responder. Em projetos maiores, vale adicionar \texttt{healthcheck}.
\end{alertbox}

\subsection{Ajustando o web-ms}

O front \texttt{web-ms} j\'a possu\'ia vari\'aveis para indicar a URL base das APIs. Antes, elas apontavam diretamente para os microsservi\c{c}os:

\begin{lstlisting}[style=code]
CLIENTES_SERVICE_URL: http://clientes-service:8081
ESTOQUE_SERVICE_URL: http://estoque-service:8082
PEDIDOS_SERVICE_URL: http://pedidos-service:8083
\end{lstlisting}

Depois, passaram a apontar para o gateway:

\begin{lstlisting}[style=code]
CLIENTES_SERVICE_URL: http://api-gateway:9000
ESTOQUE_SERVICE_URL: http://api-gateway:9000
PEDIDOS_SERVICE_URL: http://api-gateway:9000
\end{lstlisting}

\subsection{Por que o c\'odigo Java n\~ao mudou?}

O \texttt{web-ms} j\'a monta as URLs finais adicionando o caminho da API. Exemplo simplificado:

\begin{lstlisting}[style=code]
restClient.get()
    .uri(clientesServiceUrl + "/api/clientes")
    .retrieve()
    .body(ClienteResponse[].class);
\end{lstlisting}

Como \texttt{clientesServiceUrl} passou a ser \texttt{http://api-gateway:9000}, a URL final fica:

\begin{lstlisting}[style=terminal]
http://api-gateway:9000/api/clientes
\end{lstlisting}

% ================================================================
\section{Parte 5 --- Subindo e Validando a Solu\c{c}\~ao}

\subsection{Validando o Compose}

Antes de subir, \'e poss\'ivel pedir ao Docker Compose para interpretar o arquivo completo:

\begin{lstlisting}[style=terminal]
docker compose config
\end{lstlisting}

Esse comando ajuda a encontrar erros de sintaxe, indentac\~ao e configura\c{c}\~ao.


\subsection{Subindo os containers}

Depois das altera\c{c}\~oes, foi executado:

\begin{lstlisting}[style=terminal]
docker compose up -d
\end{lstlisting}

Esse comando:

\begin{itemize}[leftmargin=1.5em]
  \item baixou a imagem \texttt{nginx:1.27-alpine}, caso ela n\~ao existisse localmente;
  \item recriou containers cuja configura\c{c}\~ao mudou;
  \item iniciou o \texttt{api-gateway};
  \item manteve a aplica\c{c}\~ao rodando em segundo plano.
\end{itemize}

\subsection{Conferindo portas depois da mudan\c{c}a}

\begin{lstlisting}[style=terminal]
docker ps --format 'table {{.Names}}\t{{.Ports}}\t{{.Status}}'
\end{lstlisting}

O resultado esperado \'e semelhante a:

\begin{lstlisting}[style=terminal]
NAMES              PORTS
web-ms             0.0.0.0:9008->8080/tcp
api-gateway        0.0.0.0:9000->9000/tcp
pedidos-service    8083/tcp
clientes-service   8081/tcp
estoque-service    8082/tcp
clientes-db        5432/tcp
estoque-db         5432/tcp
pedidos-db         5432/tcp
\end{lstlisting}

\begin{infobox}[Como interpretar o resultado]
  Depois da mudan\c{c}a, apenas \texttt{api-gateway} e \texttt{web-ms} t\^em portas publicadas no host. Os servi\c{c}os \texttt{clientes-service}, \texttt{estoque-service} e \texttt{pedidos-service} continuam ativos, mas acess\'iveis apenas dentro da rede Docker.
\end{infobox}

% ================================================================
\section{Parte 6 --- Testes com curl}

\subsection{Testando clientes pelo gateway}

\begin{lstlisting}[style=terminal]
curl -s -o /tmp/clientes.out -w '%{http_code}' http://localhost:9000/api/clientes
\end{lstlisting}

Resultado esperado:

\begin{lstlisting}[style=terminal]
200
\end{lstlisting}

\subsection{Testando estoque pelo gateway}

\begin{lstlisting}[style=terminal]
curl -s -o /tmp/estoques.out -w '%{http_code}' http://localhost:9000/api/estoques
\end{lstlisting}

Resultado esperado:

\begin{lstlisting}[style=terminal]
200
\end{lstlisting}

\subsection{Testando pedidos pelo gateway}

\begin{lstlisting}[style=terminal]
curl -s -o /tmp/pedidos.out -w '%{http_code}' http://localhost:9000/api/pedidos
\end{lstlisting}

Resultado esperado:

\begin{lstlisting}[style=terminal]
200
\end{lstlisting}

\subsection{Testando o front}

\begin{lstlisting}[style=terminal]
curl -s -o /tmp/web.out -w '%{http_code}' http://localhost:9008
\end{lstlisting}

Resultado esperado:

\begin{lstlisting}[style=terminal]
200
\end{lstlisting}

\subsection{Testando uma porta antiga}

Foi testado o acesso antigo ao servi\c{c}o de clientes:

\begin{lstlisting}[style=terminal]
curl -s -o /tmp/old-clientes.out -w '%{http_code}' http://localhost:9005/api/clientes
\end{lstlisting}

Resultado observado:

\begin{lstlisting}[style=terminal]
000
\end{lstlisting}

\begin{infobox}[O que significa \texttt{000}?]
  No \texttt{curl}, o c\'odigo \texttt{000} normalmente indica que a conex\~ao n\~ao foi estabelecida. Neste caso, isso \'e esperado, pois a porta \texttt{9005} n\~ao est\'a mais publicada no host.
\end{infobox}

% ================================================================
\section{Arquitetura Final}

Depois da altera\c{c}\~ao, a arquitetura de acesso ficou assim:

\begin{center}
\begin{tabular}{lll}
\toprule
\textbf{Acesso externo} & \textbf{Gateway encaminha para} & \textbf{Responsabilidade} \\
\midrule
\texttt{localhost:9000/api/clientes} & \texttt{clientes-service:8081} & Clientes \\
\texttt{localhost:9000/api/estoques} & \texttt{estoque-service:8082} & Estoque \\
\texttt{localhost:9000/api/pedidos} & \texttt{pedidos-service:8083} & Pedidos \\
\texttt{localhost:9008} & \texttt{web-ms:8080} & Interface web \\
\bottomrule
\end{tabular}
\end{center}

\begin{alertbox}
  O \texttt{web-ms} foi mantido fora do gateway porque o objetivo desta etapa era centralizar apenas as APIs. Em outro desenho arquitetural, o Nginx tamb\'em poderia servir o front em \texttt{/} e manter as APIs em \texttt{/api}.
\end{alertbox}

% ================================================================
\section{Por que Isso Foi Feito?}

A mudan\c{c}a foi feita para aproximar o projeto de uma arquitetura mais organizada e mais comum em ambientes reais.

\subsection{Benef\'icios principais}

\begin{itemize}[leftmargin=1.5em]
  \item \textbf{Entrada \'unica}: o consumidor usa \texttt{localhost:9000} para todas as APIs.
  \item \textbf{Menos exposi\c{c}\~ao}: as portas \texttt{9005}, \texttt{9006} e \texttt{9007} deixam de ser p\'ublicas.
  \item \textbf{Roteamento centralizado}: as regras de encaminhamento ficam no Nginx.
  \item \textbf{Evolu\c{c}\~ao futura}: o gateway pode receber autentica\c{c}\~ao, HTTPS, logs, rate limit e CORS.
  \item \textbf{Cliente mais simples}: quem consome a API n\~ao precisa saber onde cada microsservi\c{c}o roda internamente.
\end{itemize}

\subsection{O que n\~ao mudou}

\begin{itemize}[leftmargin=1.5em]
  \item Os servi\c{c}os Spring Boot continuam escutando nas mesmas portas internas.
  \item Os bancos PostgreSQL continuam internos.
  \item O \texttt{pedidos-service} continua chamando \texttt{clientes-service} e \texttt{estoque-service} diretamente pela rede Docker.
  \item O front continua acess\'ivel em \texttt{localhost:9008}.
\end{itemize}


% ================================================================
\section{Resumo do Que Foi Alterado}

\begin{enumerate}[leftmargin=2em]
  \item Foi criada a pasta \texttt{nginx}.
  \item Foi criado o arquivo \texttt{nginx/api-gateway.conf}.
  \item Foi adicionado o container \texttt{api-gateway} usando \texttt{nginx:1.27-alpine}.
  \item O gateway foi publicado em \texttt{9000:9000}.
  \item Foram removidas as portas p\'ublicas dos servi\c{c}os de API.
  \item O \texttt{web-ms} passou a consumir \texttt{http://api-gateway:9000}.
  \item O Compose foi validado com \texttt{docker compose config}.
  \item Os containers foram recriados com \texttt{docker compose up -d}.
  \item Os endpoints foram testados com \texttt{curl} e retornaram \texttt{200}.
\end{enumerate}

% ================================================================
\section{Fechamento}

O ponto central desta pr\'atica \'e perceber que o API Gateway organiza a entrada do sistema sem eliminar os microsservi\c{c}os. Cada servi\c{c}o continua com sua responsabilidade, sua porta interna e seu banco quando necess\'ario. A diferen\c{c}a \'e que o consumidor externo n\~ao precisa mais conhecer todos esses detalhes.

Com o Nginx na porta \texttt{9000}, a aplica\c{c}\~ao passa a ter uma entrada \'unica para as APIs, reduz a exposi\c{c}\~ao direta dos servi\c{c}os e abre caminho para melhorias futuras como autentica\c{c}\~ao centralizada, HTTPS e monitoramento.

\end{document}
